\scanpage{37}
% archive sheet number: 36
\begingroup\makeatletter\tagsleft@true\makeatother
\uline{Corollary} \deleted{I}n the case where $\operatorname{Card}\{\ \}=4$, \inserted{and $\Sigma_1^-\ne X,\Sigma_2^-\ne X$} the only inclusion relations between the sets $\Sigma_1^+,\Sigma_1^-,\Sigma_2^+,\Sigma_2^-$ are the relations
\margin{In the form : For the \unsure{strict} relations $(*)$ to be the only inclusion relations between the sets $\Sigma_1^+,\Sigma_1^-,\Sigma_2^+,\Sigma_{\unsure{1|2}}^-$, it is necessary and sufficient that these be \inserted{\unsure{mutually|mutually}} distinct and that we have $\Sigma_1^-\ne X,\Sigma_2^-\ne X$.}
\[
\left\{\begin{aligned}
\Sigma_1^{\unsure{\top|+}}&\subset\Sigma_2^-\\
\Sigma_2^+&\subset\Sigma_1^-
\end{aligned}\right.\qquad .
\tag*{($*$)}
\]
Indeed, we do not have the reverse inclusions, by hypothesis. On the other hand, we \uline{do not have}
\[
\Sigma_1^+\subset\Sigma_1^-\quad\text{or}\quad
\Sigma_1^-\unsure{\subset|\in}\Sigma_1^+,\quad\text{for we would have}
\]
$\deleted{\Sigma_1^+,\unsure{\Sigma_1|\Sigma_1^+|\Sigma_1^-}}\ \Sigma_1^-=X$, likewise we do not have $\Sigma_2^+\subset\unsure{\deleted{\vert}}{\unsure{\Sigma|E}}_{\unsure{1|2}}^-$ nor $\Sigma_2^-\subset\Sigma_2^+$, and we do not have $\Sigma_1^+\subset\Sigma_2^+$ nor the reverse inclusion, since $\Sigma_1^+$ and $\unsure{\Sigma|\Sigma_2}^{\top}$ are disjoint and $\ne\varnothing$. It remains to see that we do not have an inclusion $\Sigma_1^-\subset\Sigma_2^-$ or reverse inclusion, we read from the diagram $(14)$ that \deleted{\unsure{this}} \inserted{the first \unsure{for example|for example}} would mean that $\deleted{\Sigma_{12}\simeq\Sigma_1^-}$ \deleted{\unsure{whence|of}} $\Sigma_2^-=X$, which we have excluded.

% The two thin, slightly wavering segments below reproduce the left marginal rule.
\noindent\begin{minipage}[t]{.04\linewidth}
\vspace{1.2\baselineskip}
\figure{\begin{tikzpicture}[x=.3pt,y=.3pt]
\draw[line width=.35pt] (0,0) -- (1,-75) -- (4,-155) -- (6,-212) -- (4,-270) -- (9,-310);
\draw[line width=.35pt] (9,-321) -- (13,-360) -- (7,-415) -- (5,-475) -- (1,-533) -- (3,-567);
\end{tikzpicture}}
\end{minipage}\hfill\begin{minipage}[t]{.94\linewidth}
\vspace{0pt}
\uline{\unsure{Definition|Definition}} We say that $\unsure{\Pi|\pi}_1,\unsure{\Pi|R}_2$ are \uline{\unsure{strictly|Strictly} non\unsure{:|.} mixed} \inserted{if they are non-mixed, and} if t\unsure{\deleted{he}}he four sets $\Sigma_1^\top,\Sigma_1^-,\Sigma_2^+,\Sigma_2^-$ are distinct, \unsure{and.|i.e.} if \unsure{((|c} in addition $\Sigma_1^-\ne X,\Sigma_2^-\ne X$ -- i.e if the only inclusion relations between the sets $\Sigma_1^\top,\Sigma_1^-,\Sigma_2^\top,\Sigma_2^-$ are $\Sigma_1^+\subset\Sigma_2^-,\Sigma_2^+\subset\Sigma_1^-$ -- which also means that we have :
\[
\Sigma_1^-\ne X,\Sigma_2^-\ne X
\quad(\text{we therefore say that }\Pi_{\unsure{1|2}},\Pi_2\text{ are \uline{proper}}),
\ E_1\ne\Sigma_{12},E_2\ne\Sigma_{12}.
\tag*{(21)}
\]
\end{minipage}
\endgroup
\clearpage
